% Uma Discussão sobre Óptica Geométrica — Cauê Takahashi de Oliveira
% Fonte reconstruída a partir do PDF original (fonte/optica-geometrica-original.pdf).
% Compilação: latexmk main.tex   (pdfLaTeX, como o original)
\documentclass[11pt,a4paper]{article}

\usepackage[T1]{fontenc}
\usepackage[utf8]{inputenc}
\usepackage{lmodern}
\usepackage[brazilian]{babel}
\usepackage[hmargin=2.7cm,vmargin=2.5cm]{geometry}
\usepackage{amsmath,amssymb,mathtools}
\usepackage{cancel}
\usepackage{microtype}
\usepackage{csquotes}
\usepackage{enumitem}
\usepackage{xcolor}
\usepackage{graphicx}
\usepackage{caption}
\usepackage{subcaption}
\usepackage[most]{tcolorbox}
\usepackage{tikz}
\usetikzlibrary{arrows.meta,decorations.markings,calc}

% ---------------------------------------------------------------- notação
\DeclareMathOperator{\tg}{tg}

% ---------------------------------------------------------------- cores
% Raios: paleta Okabe–Ito (distinguível por daltônicos).
\definecolor{azul}{HTML}{0072B2}
\definecolor{laranja}{HTML}{E69F00}
\definecolor{verde}{HTML}{009E73}
\definecolor{rosa}{HTML}{CC79A7}
\definecolor{vermelhao}{HTML}{D55E00}
\definecolor{vermelho}{HTML}{C0392B}
\definecolor{corespelho}{HTML}{3E4C66}
\definecolor{caixafundo}{HTML}{F6F8FA}
\definecolor{caixabarra}{HTML}{6283A1}
\definecolor{caixatitulo}{HTML}{1F4E79}

% ---------------------------------------------------------------- legendas
\captionsetup{font=small,labelfont={sf,bf},labelsep=period,margin=0.5cm}
\captionsetup[sub]{font=small,labelfont={rm,bf},labelsep=space,margin=0pt}

% ---------------------------------------------------------------- caixas
\newtcolorbox{comentario}{enhanced,breakable,sharp corners,boxrule=0pt,
  frame hidden,colback=caixafundo,borderline west={2.2pt}{0pt}{caixabarra},
  left=10pt,right=10pt,top=5pt,bottom=6pt,before skip=13.15pt,after skip=9.7pt,
  fontupper=\small}
\newcommand{\titulocaixa}[1]{{\sffamily\bfseries\color{caixatitulo}#1}}
\setlist[itemize]{topsep=2pt,itemsep=5pt,parsep=0pt}
\setlength{\parskip}{0.45em}

% ---------------------------------------------------------------- TikZ
\pgfdeclareradialshading{brilho}{\pgfpoint{0bp}{0bp}}{color(0bp)=(white);
  color(14bp)=(white); color(34bp)=(laranja!55); color(50bp)=(white)}
\tikzset{
  >={Stealth[length=6.2pt,width=5.1pt]},
  % raio de luz com seta no meio do segmento
  raio/.style={line width=0.9pt,postaction={decorate},
    decoration={markings,mark=at position 0.5 with {\arrow{>}}}},
  eixo/.style={line width=0.4pt,black!50},
  espelho/.style={line width=1.6pt,corespelho},
  hachura/.style={line width=0.4pt,corespelho!60},
  ponto/.style={circle,fill,inner sep=0pt,minimum size=3pt},
  cota/.style={line width=0.5pt,shorten <=1pt,shorten >=1pt,
    {Latex[length=3.3pt,width=2.5pt]}-{Latex[length=3.3pt,width=2.5pt]}},
  guia/.style={line width=0.4pt,black!50,dotted},
  anteparo/.style={line width=0.6pt,black!65,dashed,fill=black!4},
  % brilho em torno de uma fonte pontual
  brilho/.pic={\shade[shading=brilho] (0,0) circle (0.21); \fill (0,0) circle (0.06);},
  % olho voltado para a esquerda (use xscale=-1 para a direita)
  olho/.pic={
    \draw[line width=0.6pt,fill=white] (0,0) circle (0.345);
    \draw[line width=0.6pt,fill=white] (150:0.345) arc[start angle=131,end angle=229,radius=0.224];
    \fill[azul!70!black] (-0.316,0) ellipse[x radius=0.034,y radius=0.138];
    \draw[line width=0.4pt,fill=black!8] (-0.219,0) ellipse[x radius=0.052,y radius=0.121];
    \draw[line width=0.6pt] (0.345,0) -- (0.483,0);},
}

\usepackage[hidelinks,pdftitle={Uma Discussão sobre Óptica Geométrica},
  pdfauthor={Cauê Takahashi de Oliveira}]{hyperref}

\title{\textbf{Uma Discussão sobre Óptica Geométrica}\\[11pt]
  {\large\itshape Um pequeno trabalho}}
\author{Cauê Takahashi de Oliveira}
\date{Novembro de 2026}

\begin{document}
\maketitle
\thispagestyle{empty}

Este é um trabalho autoral, em que vamos discutir de onde vêm as regras que
aprendemos para espelhos e lentes. Em vez de aceitar essas regras como prontas,
vamos entender o que significa, fisicamente, formar uma imagem, por que dizemos
que ela se forma no cruzamento de feixes luminosos e por que o estudo dos
espelhos côncavos usa apenas alguns feixes convenientes. A partir disso,
demonstraremos que todo feixe que chega paralelo ao eixo cruza o mesmo ponto, a
uma distância $R/2$ do vértice, o que dá origem ao conceito de foco, e seguiremos
rumo à equação de Gauss.

Nesta seção vamos compreender e demonstrar o porquê de as fórmulas de óptica
fazerem sentido e o que de fato os modelos teóricos dizem sobre a realidade.

\medskip
\noindent\textbf{1º)} O problema do estudo da óptica começa cedo: é normal que
alunos compreendam bem os princípios da luz, reflexões e espelhos planos, mas,
ao estudar lentes e espelhos côncavos, se percam completamente, sem que consigam
estabelecer uma conexão direta entre os fundamentos estudados e essa nova
frente, parecendo até que algumas regras foram inventadas.

\medskip
\noindent\textbf{2º)} A primeira objeção surge logo no início, onde é jogado o
fato de que no cruzamento de raios luminosos uma imagem é formada. A lógica por
trás, por incrível que pareça, é bem simples. Ninguém discorda ou duvida que o
fenômeno da visão ocorre quando feixes de luz provenientes de um objeto,
carregando informações, atravessam nosso olho e são interpretados pelo nosso
cérebro. Nesse ponto, ainda é óbvio para todos que um corpo emite infinitos
feixes de luz (Figura~\ref{fig:ponto}). Entretanto, ao estudar espelhos
côncavos, utilizam-se apenas dois feixes, os quais ainda por cima seguem regras
estranhas: \enquote{o feixe que chega paralelo passa pelo foco}
(Figura~\ref{fig:dois-feixes}). Por quê?

\begin{figure}[htbp]
  \centering
  \begin{subfigure}[b]{0.3\textwidth}
    \centering
    % Figura 1a: um ponto do objeto emite feixes em todas as direções.
\begin{tikzpicture}[font=\small]
  \foreach \a in {0,15,...,345}
    \draw[line width=0.9pt,laranja!85!black,postaction={decorate},
      decoration={markings,mark=at position 0.74 with {\arrow{>}}}]
      (0,0) -- (\a:1.55);
  \foreach \a in {7.5,22.5,...,352.5}
    \draw[line width=0.4pt,laranja!60,postaction={decorate},
      decoration={markings,mark=at position 0.8 with {\arrow{>}}}]
      (0,0) -- (\a:1.2);
  \pic at (0,0) {brilho};
  \node[anchor=base] at (0,-2.0) {um ponto do objeto};
\end{tikzpicture}

    \caption{}\label{fig:ponto}
  \end{subfigure}%
  \hfill
  \begin{subfigure}[b]{0.66\textwidth}
    \centering
    % Figura 1b: as duas regras usuais para o espelho côncavo.
% Vértice V na origem; R = 4, C = (-R, 0), F = (-R/2, 0). Objeto em p = 6.
\begin{tikzpicture}[font=\small]
  \coordinate (C) at (-4,0);
  \coordinate (F) at (-2,0);
  \draw[eixo] (-6.6,0) -- (0.6,0);
  \draw[espelho] (C) ++(-22:4) arc[start angle=-22,end angle=22,radius=4];
  \foreach \a in {-20,-17.5,...,20}
    \draw[hachura] (C) ++(\a:4) -- ++(\a-40:0.18);
  % objeto e imagem (p' = 3, altura -0,4)
  \draw[line width=1.1pt,-{Stealth[length=3.5pt,width=2.8pt]}] (-6,0) -- (-6,0.8);
  \draw[line width=1.1pt,black!50,-{Stealth[length=3.5pt,width=2.8pt]}] (-3,0) -- (-3,-0.4);
  % (1) chega paralelo -> passa pelo foco
  \draw[raio,vermelho] (-6,0.8) -- (-0.081,0.8);
  \draw[raio,vermelho] (-0.081,0.8) -- (-3.705,-0.711);
  % (2) passa pelo foco -> sai paralelo
  \draw[raio,azul] (-6,0.8) -- (-0.02,-0.4);
  \draw[raio,azul] (-0.02,-0.4) -- (-3.705,-0.4);
  \fill (C) circle (0.053) (F) circle (0.053);
  \node[below=1pt] at (C) {$C$};
  \node[below=1pt] at (F) {$F$};
  \node[anchor=north west,inner sep=1pt] at (0.12,-0.12) {$V$};
  \node[anchor=base,font=\scriptsize,text=vermelho] at (-3.17,1.025)
    {(1) chega paralelo $\to$ passa pelo foco};
  \node[anchor=base,font=\scriptsize,text=azul] at (-2.98,-1.297)
    {(2) passa pelo foco $\to$ sai paralelo};
\end{tikzpicture}

    \caption{}\label{fig:dois-feixes}
  \end{subfigure}
  \caption{(a) Cada ponto de um corpo emite infinitos feixes de luz, em todas as
    direções. (b) No estudo dos espelhos côncavos, porém, costumam-se usar apenas
    dois feixes, com regras como \enquote{o feixe que chega paralelo passa pelo
    foco} (1) e \enquote{o feixe que passa pelo foco sai paralelo} (2).
    $C$: centro de curvatura; $F$: foco; $V$: vértice.}
  \label{fig:raios}
\end{figure}

\medskip
\noindent\textbf{3º)} Veja bem, antes de discutirmos e mostrarmos o porquê,
vamos entender o que significa formar uma imagem. O que acontece de fato é que a
informação que a luz carrega, embora não seja perdida, é rearranjada por um
aparato refletor. Instantaneamente, o rearranjo gerado impede que nós
interpretemos o objeto. Suponha o seguinte exemplo: duas pessoas chegam em você
e dizem: \enquote{Olá, tudo bem com você?} e a outra: \enquote{Bem você? Com,
olá tudo}. Note que, embora o conjunto de palavras seja o mesmo e cada palavra
carregue a sua mesma informação própria, ainda assim sua interpretação foi
prejudicada, apenas rearranjando informação (Figura~\ref{fig:palavras}).
Analogamente, cada feixe de luz carrega uma informação única e, ao
interpretarmos esse conjunto de dados na ordem imposta pela emissão do objeto,
interpretamos ele (Figura~\ref{fig:feixes-ordem}).

A ideia de formar uma imagem surge fundamentalmente desse conceito; veja como
tudo se encaixa: a luz sai de um objeto, a luz reflete em um espelho côncavo e
cada feixe de luz tem sua trajetória redefinida. Se eles saírem de um ponto, por
exemplo, essa emissão pontual ordenada configura nossa interpretação do objeto e
nossa distância dele. É óbvio que, caso a luz que foi espalhada seja
reconfigurada novamente de forma que todos os feixes luminosos saiam na mesma
configuração de emissão inicial, nesse caso um ponto, voltaríamos a interpretar
aquele conjunto de informação da mesma maneira que o objeto. Por isso, dizer que
a imagem se forma no encontro de feixes luminosos não significa que magicamente
naquele ponto se formará um holograma do objeto, e sim que, ao se cruzarem,
assumem a mesma configuração da informação emitida pelo corpo inicialmente e,
portanto, caso um anteparo seja colocado lá, a projeção vista seria interpretada
da mesma maneira que interpretamos o objeto (Figura~\ref{fig:imagem-real}).
Denominamos esse fenômeno de formação de imagem. Note que nem sempre será
possível colocar um anteparo no local em que a imagem se forma, e isso não muda o
fato de que a imagem se formará lá (Figura~\ref{fig:espelho-plano}).

\begin{figure}[htbp]
  \centering
  \begin{subfigure}{\textwidth}
    \centering
    % Figura 2a: as mesmas palavras em outra ordem.
\begin{tikzpicture}[font=\small,
    palavra/.style={draw=#1,fill=#1!12,line width=0.7pt,rounded corners=2.5pt,
      minimum width=1.1cm,minimum height=0.6cm,inner sep=1pt,font=\sffamily\small}]
  \node[inner sep=0pt,anchor=base east] at (2.75,1.577) {Pessoa 1:};
  \node[inner sep=0pt,anchor=base east] at (2.75,0.627) {Pessoa 2:};
  \foreach \w/\c [count=\i from 0] in {{Olá,}/azul,tudo/laranja,bem/verde,com/rosa,você?/vermelhao}
    \node[palavra=\c] at (3.657+1.275*\i,1.698) {\w};
  \foreach \w/\c [count=\i from 0] in {Bem/verde,você?/vermelhao,{Com,}/rosa,olá/azul,tudo/laranja}
    \node[palavra=\c] at (3.657+1.275*\i,0.748) {\w};
  \node[inner sep=0pt,anchor=base west,text=verde!70!black] at (9.953,1.61)
    {$\checkmark$ entendemos a mensagem};
  \node[inner sep=0pt,anchor=base west,text=vermelhao!80!black] at (9.953,0.661)
    {$\times$ interpretação prejudicada};
\end{tikzpicture}

    \caption{Mesmas palavras, cada uma com a sua informação; só a ordem mudou.}
    \label{fig:palavras}
  \end{subfigure}

  \bigskip
  \begin{subfigure}{\textwidth}
    \centering
    % Figura 2b: os mesmos feixes, na ordem da emissão e rearranjados.
\begin{tikzpicture}[font=\small]
  % à esquerda: na ordem da emissão
  \foreach \dy/\c in {0.3/azul,0.15/laranja,0/verde,-0.15/rosa,-0.3/vermelhao}
    \draw[raio,\c] (2.162,2.181) -- (5.412,2.181+\dy);
  \pic at (2.162,2.181) {brilho};
  \pic at (5.782,2.181) {olho};
  \node[anchor=base,align=center] at (4.021,0.634)
    {na ordem da emissão:\\o olho interpreta um ponto};
  % à direita: rearranjados por um refletor irregular
  \fill[black!9] (8.762,1.431) -- (9.012,1.541) -- (9.362,1.371) -- (9.612,1.561)
    -- (9.862,1.391) -- (10.112,1.541) -- (10.362,1.351) -- (10.612,1.521)
    -- (10.862,1.381) -- (11.112,1.551) -- (11.362,1.371) -- (11.612,1.511)
    -- (11.862,1.431) -- (11.862,1.131) -- (8.762,1.131) -- cycle;
  \draw[line width=1.2pt,corespelho] (8.762,1.431) -- (9.012,1.541) -- (9.362,1.371)
    -- (9.612,1.561) -- (9.862,1.391) -- (10.112,1.541) -- (10.362,1.351)
    -- (10.612,1.521) -- (10.862,1.381) -- (11.112,1.551) -- (11.362,1.371)
    -- (11.612,1.511) -- (11.862,1.431);
  \foreach \px/\py/\qy/\c in {9.362/1.371/2.751/azul,9.862/1.391/3.051/laranja,
      10.362/1.351/3.201/verde,10.862/1.381/2.901/rosa,11.362/1.371/3.351/vermelhao} {
    \draw[raio,\c] (8.562,2.981) -- (\px,\py);
    \draw[raio,\c] (\px,\py) -- (13.012,\qy);
  }
  \pic at (8.562,2.981) {brilho};
  \pic at (13.382,3.051) {olho};
  \node[anchor=base west,text=vermelhao!80!black,font=\large\bfseries] at (13.16,3.537) {?};
  \node[anchor=base,align=center] at (10.821,0.384)
    {rearranjados por um refletor irregular:\\a mesma informação, fora de ordem};
\end{tikzpicture}

    \caption{O mesmo vale para os feixes de luz.}
    \label{fig:feixes-ordem}
  \end{subfigure}
  \caption{A analogia das palavras. As cores indicam a informação que cada
    palavra (ou cada feixe) carrega. Em (a), as duas frases têm exatamente as
    mesmas palavras, mas só a primeira é interpretada. Em (b), na ordem imposta
    pela emissão, o olho interpreta os feixes como vindos de um ponto; os mesmos
    feixes, rearranjados por uma superfície irregular, chegam fora de ordem e o
    ponto não é reconhecido.}
  \label{fig:analogia}
\end{figure}

\begin{figure}[htbp]
  \centering
  % Figura 3: imagem real em um espelho côncavo.
% C na origem, R = 7,84; objeto A a p = 1,5R do vértice, imagem A' a p' = 0,75R
% (1/p + 1/p' = 2/R), altura -0,5 da do objeto. Os raios são desenhados de forma
% idealizada (todos passam por A').
\begin{tikzpicture}[font=\small]
  \def\R{7.84}
  \coordinate (C) at (0,0);
  \coordinate (F) at (\R/2,0);
  \coordinate (A) at (-\R/2,\R/5);
  \coordinate (A') at (\R/4,-\R/10);
  \draw[eixo] (-4.41,0) -- (8.428,0);
  \draw[espelho] (-12.5:\R) arc[start angle=-12.5,end angle=12.5,radius=\R];
  \foreach \a in {-12,-10.5,...,12}
    \draw[hachura] (\a:\R) -- ++(\a-40:0.195);
  \foreach \y/\c in {0.15/azul,0.075/laranja,0/verde,-0.075/rosa,-0.15/vermelhao} {
    \pgfmathsetmacro{\yy}{\y*\R}
    \pgfmathsetmacro{\xx}{sqrt(\R*\R-\yy*\yy)}
    \coordinate (M) at (\xx,\yy);
    \draw[raio,\c] (A) -- (M);
    \draw[raio,\c] (M) -- (A');
    \draw[raio,\c] (A') -- ($(A')!-0.29!(M)$);
  }
  \fill (C) circle (0.053) (F) circle (0.053);
  \node[fill=white,inner sep=1pt] at (0,0.267) {$C$};
  \draw[line width=0.3pt,black!50] (\R/2,0.049) -- (\R/2,0.588);
  \node[fill=white,inner sep=1pt] at (\R/2,0.776) {$F$};
  \node[anchor=north west,inner sep=1pt] at (\R+0.14,-0.07) {$V$};
  \draw[anteparo] (1.891,-1.715) rectangle (2.029,-0.117);
  \node[anchor=base,text=black!80] at (1.99,-2.101) {anteparo};
  \pic at (A) {brilho};
  \fill (A') circle (0.059);
  \node[above left=0pt] at (A) {$A$};
  \node[above left=-1pt] at (A') {$A'$};
  \pic[xscale=-1,scale=1.022] at (-0.078,-1.009) {olho};
\end{tikzpicture}

  \caption{Os feixes que saem do ponto $A$ têm sua trajetória redefinida pelo
    espelho côncavo e voltam a se encontrar em $A'$, onde assumem de novo a
    configuração de uma emissão pontual (as cores mostram que nenhuma informação
    se perde). Um anteparo colocado em $A'$ mostraria um ponto luminoso; sem o
    anteparo, os feixes seguem e chegam ao olho como se tivessem sido emitidos
    em $A'$.}
  \label{fig:imagem-real}
\end{figure}

\begin{figure}[htbp]
  \centering
  % Figura 4: imagem virtual no espelho plano (A' simétrico de A).
\begin{tikzpicture}[font=\small]
  \fill[black!7] (6.613,-1.558) rectangle (6.833,2.442);
  \draw[espelho] (6.613,-1.558) -- (6.613,2.442);
  \foreach \y in {-1.458,-1.208,...,2.3}
    \draw[hachura] (6.613,\y) -- ++(0.2,0.12);
  \foreach \yp/\yo/\c in {0.242/-0.358/azul,0.042/-0.678/verde,-0.158/-0.998/vermelhao} {
    \draw[raio,\c] (3.613,1.242) -- (6.613,\yp);
    \draw[raio,\c] (6.613,\yp) -- (4.813,\yo);
    \draw[line width=0.6pt,dashed,\c] (6.613,\yp) -- (9.613,1.242);
  }
  \pic at (3.613,1.242) {brilho};
  \draw[line width=0.6pt,fill=white] (9.613,1.242) circle (0.077);
  \node[anchor=base] at (6.615,2.695) {espelho plano};
  \node[above left=0pt] at (3.613,1.242) {$A$};
  \node[above left=0pt] at (9.613,1.242) {$A'$};
  \draw[anteparo] (10.013,0.492) rectangle (10.153,1.992);
  \node[inner sep=0pt,anchor=base west,align=left,text=black!85] at (10.448,0.743)
    {anteparo:\\nenhuma luz\\chega aqui};
  \pic[xscale=-1,scale=0.957] at (4.473,-0.678) {olho};
\end{tikzpicture}

  \caption{Nem sempre é possível colocar um anteparo onde a imagem se forma. No
    espelho plano, os feixes refletidos chegam ao olho como se saíssem de $A'$,
    atrás do espelho (prolongamentos tracejados): a imagem se forma lá, mas um
    anteparo em $A'$ não receberia luz alguma.}
  \label{fig:espelho-plano}
\end{figure}

\clearpage
\section*{Demonstrando certas propriedades e a equação de Gauss}

\begin{itemize}[label=$\ast$,leftmargin=17pt,labelsep=6pt,topsep=4pt]
  \item Na vida real, um espelho côncavo possui grau de abertura, no geral,
    menor ou igual a $10^\circ$. Como essa é a abertura total, do vértice até o
    topo do espelho o ângulo é só a metade, $\theta \le 5^\circ$; logo, a altura do
    vértice do espelho até seu topo, $H = R\sin\theta \le R\sin 5^\circ \approx
    0{,}09\,R$, é muito menor que o raio, e qualquer ponto do espelho, a uma
    altura $h$ do eixo, tem $h \le H \ll R$ (Figura~\ref{fig:abertura}).
\end{itemize}

\begin{figure}[htbp]
  \centering
  % Figura 5: abertura real (em escala). C na origem, R = 6,5 cm.
% Abertura total phi = 10°: o espelho vai de -5° a 5°; do vértice ao topo, theta = 5°.
\begin{tikzpicture}[font=\small]
  \def\R{6.5}
  \draw[eixo] (-0.4,0) -- (7.6,0);
  \draw[line width=0.5pt,black!30,dashed] (5:\R) arc[start angle=5,end angle=12,radius=\R];
  \draw[line width=0.5pt,black!30,dashed] (-5:\R) arc[start angle=-5,end angle=-12,radius=\R];
  \draw[espelho] (-5:\R) arc[start angle=-5,end angle=5,radius=\R];
  \foreach \a in {-4.5,-3,...,4.5}
    \draw[hachura] (\a:\R) -- ++(\a-40:0.18);
  \draw[line width=0.6pt] (0,0) -- (5:\R);
  \draw[line width=0.4pt] (0:4.6) arc[start angle=0,end angle=5,radius=4.6];
  \node[anchor=base] at (3.0,0.42) {$R$};
  \node[inner sep=0pt,anchor=base west] at (4.75,0.1) {$\theta$};
  \draw[guia] (5:\R) -- (7.35,{\R*sin(5)});
  \draw[cota] (7.2,0) -- (7.2,{\R*sin(5)});
  \node[inner sep=0pt,anchor=base west] at (7.337,0.2) {$H$};
  \fill (0,0) circle (0.053);
  \node[anchor=base] at (-0.0,-0.435) {$C$};
  \node[inner sep=0pt,anchor=base west] at (6.656,-0.395) {$V$};
  \node[anchor=base] at (3.4,-0.988)
    {$\varphi \le 10^\circ \implies \theta = \varphi/2 \le 5^\circ \implies H = R\sin\theta \lesssim 0{,}09\,R$};
\end{tikzpicture}

  \caption{Abertura de um espelho côncavo real, desenhada em escala. Com
    abertura total $\varphi \le 10^\circ$ (ou seja, $\theta = \varphi/2 \le 5^\circ$
    do vértice até o topo), a altura $H$ do vértice até o topo do espelho não
    passa de cerca de 9\% do raio $R$, e o espelho é quase plano perto do
    vértice.}
  \label{fig:abertura}
\end{figure}

Queremos resolver o seguinte problema: dadas as dimensões de um espelho côncavo
e a posição de um corpo, onde se formará sua imagem? Isto é, para um corpo
puntiforme, o ponto de convergência de todos os seus feixes. Note que, como a
luz se propaga em linha reta, dois feixes podem se intersectar em somente um
ponto e, assumindo que todos os feixes devem se cruzar neste ponto, basta
equacionar o cruzamento de dois feixes distintos e poderemos afirmar que no
ponto em que eles se cruzarem todos os feixes deverão se cruzar. Assim, cabe a
nós escolhermos feixes convenientes para equacionar
(Figura~\ref{fig:demonstracao}).

\begin{figure}[htbp]
  \centering
  % Figura 6: a figura da demonstração (dimensões exageradas).
% C na origem; R = 6,48; O = (R - p, h) com h = 1,729 e p = 4,536.
% sin(theta) = h/R; o feixe delta faz o ângulo 2*theta com o eixo;
% o feixe O -> V reflete-se com o mesmo ângulo alpha, tg(alpha) = h/p.
\begin{tikzpicture}[font=\small]
  \def\R{6.48}
  \def\h{1.729}
  \def\p{4.536}
  \pgfmathsetmacro{\th}{asin(\h/\R)}
  \pgfmathsetmacro{\al}{atan(\h/\p)}
  \pgfmathsetmacro{\xP}{sqrt(\R*\R-\h*\h)}
  \coordinate (C) at (0,0);
  \coordinate (O) at ({\R-\p},\h);
  \coordinate (P) at (\xP,\h);
  \coordinate (V) at (\R,0);
  % eixos
  \draw[line width=0.5pt,-{Stealth[length=4pt,width=3pt]}] (-2.916,0) -- (10.235,0);
  \draw[line width=0.5pt,-{Stealth[length=4pt,width=3pt]}] (0,-3.996) -- (0,4.403);
  \node[inner sep=0pt,anchor=base west] at (9.922,-0.288) {$x$};
  \node[inner sep=0pt,anchor=base west] at (-0.322,4.14) {$y$};
  % espelho e altura H
  \draw[espelho] (-34:\R) arc[start angle=-34,end angle=34,radius=\R];
  \foreach \a in {-32.5,-30,...,32.5}
    \draw[hachura] (\a:\R) -- ++(\a-40:0.215);
  \draw[guia] (34:\R) -- (9.774,{\R*sin(34)});
  \draw[cota] (9.612,0) -- (9.612,{\R*sin(34)});
  \node[inner sep=0pt,anchor=base west] at (9.749,1.692) {$H$};
  % normal n (passa por C) e tangente r em P
  \draw[line width=0.6pt,black!65,dashed] (C) -- (\th:8.037);
  \node[inner sep=0pt,anchor=base west] at (7.885,2.37) {$n$ (normal)};
  \draw[line width=0.6pt,black!65] ($(P)+(\th-90:1.296)$) -- ($(P)+(\th+90:1.296)$);
  \node[inner sep=0pt,anchor=base west] at (6.041,2.89) {$r$ (tangente)};
  \draw[line width=0.4pt] ($(P)+(\th+90:0.22)$) -- ++(\th+180:0.22) -- ++(\th-90:0.22);
  % cotas p e R
  \draw[line width=0.5pt,dashed] (O) -- (O|-C);
  \draw[guia] (O|-C) -- ++(0,-0.777);
  \draw[guia] (V) -- ++(0,-1.533);
  \draw[cota] ({\R-\p},-0.594) -- (\R,-0.594)
    node[midway,fill=white,inner sep=1.5pt] {$p$};
  \draw[cota] (0,-1.35) -- (\R,-1.35) node[midway,fill=white,inner sep=1.5pt] {$R$};
  % feixe paralelo e feixe refletido delta
  \draw[raio,vermelho] (O) -- (P);
  \draw[raio,vermelho] (P) -- ++(180+2*\th:2.086) coordinate (D);
  \draw[line width=0.8pt,vermelho,densely dotted] (D) -- ++(180+2*\th:0.773);
  \node[inner sep=0pt,anchor=base west,text=vermelho,font=\normalsize\bfseries] at (3.397,0.179) {?};
  \node[inner sep=0pt,anchor=base west,text=vermelho] at (5.6,0.722) {$\delta$};
  % feixe O -> V e seu reflexo
  \draw[raio,azul] (O) -- (V);
  \draw[raio,azul] (V) -- ++(180+\al:4.123);
  % ângulos
  \draw[line width=0.4pt] (1,0) arc[start angle=0,end angle=\th,radius=1];
  \node[inner sep=0pt,anchor=base west] at (1.181,0.05) {$\theta$};
  \draw[line width=0.4pt] ($(P)+(180:1.202)$) arc[start angle=180,end angle=180+\th,radius=1.202];
  \node[inner sep=0pt,anchor=base west] at ($(P)+(-1.576,-0.325)$) {$\theta$};
  \draw[line width=0.4pt] ($(P)+(180+\th:1.5)$) arc[start angle=180+\th,end angle=180+2*\th,radius=1.5];
  \node[inner sep=0pt,anchor=base west] at ($(P)+(-1.716,-0.82)$) {$\theta$};
  \draw[line width=0.4pt,azul] ($(V)+(180-\al:1.301)$) arc[start angle=180-\al,end angle=180+\al,radius=1.301];
  \node[inner sep=0pt,anchor=base west,text=azul] at ($(V)+(-1.674,0.211)$) {$\alpha$};
  \node[inner sep=0pt,anchor=base west,text=azul] at ($(V)+(-1.674,-0.362)$) {$\alpha$};
  % pontos e rótulos
  \fill (O) circle (0.07) (P) circle (0.07) (V) circle (0.07);
  \node[inner sep=0pt,anchor=base west] at (0.974,2.029) {$O\ (R-p,\ h)$};
  \node[inner sep=0pt,anchor=base west] at (6.727,1.468) {$P = (x_P,\ h)$};
  \node[inner sep=0pt,anchor=base west] at (6.681,-0.442) {$V = (R,\ 0)$};
  \node[inner sep=0pt,anchor=base west] at (-0.412,-0.375) {$C$};
\end{tikzpicture}

  \caption{Redesenho da figura da demonstração (dimensões exageradas). O centro
    $C$ do espelho está na origem. O ponto $O$ emite um feixe paralelo ao eixo,
    que atinge o espelho em $P$ e é refletido como o feixe $\delta$; onde
    $\delta$ cruza o eixo das abscissas é justamente o que vamos calcular. Um
    segundo feixe, de $O$ até o vértice $V$, reflete-se simetricamente em relação
    ao eixo (ângulos $\alpha$). $r$: reta tangente ao espelho em $P$; $n$: reta
    normal; $H$: altura do espelho.}
  \label{fig:demonstracao}
\end{figure}

\noindent\textbf{i)} Primeiro calcula-se onde o raio paralelo deve cruzar o eixo
das abscissas. Para isso, devemos equacionar a reta normal, para achar $\theta$;
depois de fazermos isso, equacionamos o feixe $\delta$ e descobrimos o ponto de
cruzamento.

\begin{comentario}
  \titulocaixa{Comentário}\quad O feixe que sai de $O$ paralelo ao eixo atinge o
  espelho no ponto $P$, à mesma altura $h$. A reta normal $n$ ao espelho em $P$
  passa pelo centro $C$ (o raio de uma circunferência é sempre perpendicular a
  ela), e a reta tangente $r$ é perpendicular a $n$. Como o feixe incidente é
  paralelo ao eixo $x$, o ângulo de incidência em $P$ (entre o feixe e a normal)
  é igual ao ângulo $\theta$ que a normal faz com o eixo em $C$ (ângulos
  alternos internos). Pela lei da reflexão, o feixe refletido $\delta$ faz o
  mesmo ângulo $\theta$ do outro lado da normal; por isso, $\delta$ faz com o
  eixo $x$ um ângulo $2\theta$. Conhecer $\theta$ é, portanto, o primeiro passo
  para equacionar $\delta$.
\end{comentario}

Equação da circunferência (centro $C$ na origem e raio $R$):
\begin{equation}\label{eq:circ}
  x^2 + y^2 = R^2 \quad\Longrightarrow\quad 2x + 2y\frac{dy}{dx} = 0
  \quad\Longrightarrow\quad \frac{dy}{dx} = -\frac{x}{y}.
\end{equation}
Como $r \perp n$:
\begin{equation}\label{eq:normal}
  m = -\left(\frac{dy}{dx}\right)^{-1} \quad\Longrightarrow\quad m = \frac{y}{x}.
\end{equation}
Para $y = h$ (e, para o espelho, escolhe-se o sinal $+$):
\begin{equation}\label{eq:xP}
  x^2 + h^2 = R^2 \quad\Longrightarrow\quad x = \pm\sqrt{R^2 - h^2}
  \quad\Longrightarrow\quad x_P = \sqrt{R^2 - h^2}.
\end{equation}
Logo,
\begin{equation}\label{eq:m}
  m = \frac{h}{\sqrt{R^2 - h^2}} = \frac{1}{\sqrt{\dfrac{R^2}{h^2} - 1}} = \tg\theta.
\end{equation}
Como $h \ll R \Rightarrow \dfrac{R}{h} \gg 1$, logo
\begin{equation}\label{eq:aprox}
  m \approx \frac{1}{\sqrt{\dfrac{R^2}{h^2}}} = \frac{h}{R}
  \qquad\Longrightarrow\qquad \tg\theta \approx \frac{h}{R}.
\end{equation}
Essa aproximação é muito boa. Considerando um espelho médio, de $8^\circ$ de
abertura, temos $H = R\sin 4^\circ$, e $\sin 4^\circ \approx 0{,}07$; elevando ao
quadrado,
\[
  \frac{R^2}{H^2} \approx \frac{1}{(0{,}07)^2} \approx 204,
\]
um valor acima de 200 (e, como $h \le H$, $R^2/h^2$ é maior ainda). Perto de 200,
o $1$ se torna realmente muito pequeno, menos de 0,5\%, e não faz diferença
nenhuma. Principalmente porque depois ainda passamos uma raiz: sair de
$\sqrt{204}$ para $\sqrt{203}$ é um erro ainda menor, de cerca de 0,25\%.

\begin{comentario}
  \titulocaixa{Explicando as passagens}
  \begin{itemize}[leftmargin=12pt,labelsep=6pt]
    \item Em \eqref{eq:circ}, derivamos os dois lados da equação em relação a
      $x$, tratando $y$ como função de $x$ (derivação implícita): pela regra da
      cadeia, $\frac{d}{dx}(y^2) = 2y\frac{dy}{dx}$, e a derivada de $R^2$, que é
      constante, é zero. O resultado $\frac{dy}{dx}$ é a inclinação da reta
      tangente $r$ no ponto $(x, y)$ do espelho.
    \item Em \eqref{eq:normal}, usamos que duas retas perpendiculares têm
      inclinações cujo produto é $-1$. Repare que $m = y/x$ é exatamente a
      inclinação da reta que liga a origem $C$ ao ponto $(x, y)$: a conta
      confirma que a normal passa pelo centro do espelho.
    \item Em \eqref{eq:xP}, a reta $y = h$ corta a circunferência em dois pontos;
      o espelho está do lado positivo do eixo $x$, por isso fica a raiz positiva
      e $P = (x_P, h)$.
    \item Em \eqref{eq:m}, a inclinação de uma reta é a tangente do ângulo que
      ela faz com o eixo $x$; como a normal faz o ângulo $\theta$ com o eixo,
      $m = \tg\theta$. A segunda forma vem de dividir numerador e denominador
      por $h$.
    \item Em \eqref{eq:aprox}, como $R/h \gg 1$, o termo $-1$ dentro da raiz é
      desprezível diante de $R^2/h^2$. Por exemplo, para $h = 0{,}1\,R$: o valor
      exato é $\tg\theta = 1/\sqrt{99} \approx 0{,}1005$ e a aproximação dá
      $h/R = 0{,}1000$, um erro de apenas 0,5\%.
  \end{itemize}
\end{comentario}

Como $m = \tg\theta \approx \dfrac{h}{R}$:
\begin{equation}\label{eq:tg2}
  \tg 2\theta = \frac{2m}{1 - m^2} = \frac{2h}{R}\cdot\frac{1}{1 - \dfrac{h^2}{R^2}}
  \quad\Longrightarrow\quad \tg 2\theta = \frac{2hR}{R^2 - h^2},
\end{equation}
logo o feixe $\delta$: $y - \tg 2\theta\, x - b = 0$. Assim,
\begin{equation}\label{eq:b}
  b = h - \tg 2\theta\,\sqrt{R^2 - h^2}.
\end{equation}
Logo,
\begin{equation}\label{eq:delta}
  \delta:\quad y - \tg 2\theta\, x - h + \tg 2\theta\,\sqrt{R^2 - h^2} = 0.
\end{equation}
Tomando $y = 0 \Rightarrow x = x_0$ ($x$ associado ao cruzamento do eixo $x$):
\begin{align}
  0 - \tg 2\theta\, x_0 - h + \tg 2\theta\,\sqrt{R^2 - h^2} = 0
  \quad&\Longrightarrow\quad x_0 = \sqrt{R^2 - h^2} - \frac{h}{\tg 2\theta}
  \label{eq:x0}\\
  &\Longrightarrow\quad x_0 = \sqrt{R^2 - h^2} - \frac{h\,(R^2 - h^2)}{2hR}
  \label{eq:x0b}\\
  &\Longrightarrow\quad x_0 = \sqrt{R^2 - h^2} - \frac{R}{2} + \frac{h^2}{2R}.
  \notag
\end{align}
Dividindo por $R$:
\begin{equation}\label{eq:x0R}
  \frac{x_0}{R} = \sqrt{1 - \frac{h^2}{R^2}} - \frac{1}{2}
  + \frac{1}{2}\left(\frac{h}{R}\right)^{2}, \qquad \text{com}\ \frac{h}{R} \to 0.
\end{equation}
Logo,
\begin{equation}\label{eq:foco}
  \frac{x_0}{R} = 1 - \frac{1}{2} = \frac{1}{2}
  \quad\Longrightarrow\quad \boxed{x_0 = \frac{R}{2}}
\end{equation}

\begin{comentario}
  \titulocaixa{Explicando as passagens}
  \begin{itemize}[leftmargin=12pt,labelsep=6pt]
    \item Em \eqref{eq:tg2}, usamos a fórmula da tangente do arco duplo,
      $\tg 2\theta = \dfrac{2\tg\theta}{1 - \tg^2\theta}$, com
      $\tg\theta = m \approx h/R$; depois multiplicamos numerador e denominador
      por $R^2$.
    \item Em \eqref{eq:b} e \eqref{eq:delta}, o feixe $\delta$ é uma reta de
      coeficiente angular $\tg 2\theta$ (ela faz o ângulo $2\theta$ com o eixo)
      que passa por $P = \bigl(\sqrt{R^2 - h^2},\, h\bigr)$; substituindo as
      coordenadas de $P$ na equação da reta, encontramos o coeficiente linear
      $b$.
    \item Em \eqref{eq:x0}, o cruzamento com o eixo das abscissas é o ponto da
      reta onde $y = 0$; basta isolar $x_0$.
    \item Em \eqref{eq:x0b}, como $\dfrac{1}{\tg 2\theta} = \dfrac{R^2 - h^2}{2hR}$,
      o $h$ se cancela e sobra $\dfrac{R^2 - h^2}{2R} = \dfrac{R}{2} - \dfrac{h^2}{2R}$.
    \item Em \eqref{eq:x0R} e \eqref{eq:foco}, como $h \ll R$, a razão $h/R$ é
      praticamente zero: a raiz vai para 1 e o termo $\frac{1}{2}(h/R)^2$
      desaparece. O resultado \emph{não depende de $h$}: qualquer feixe que chega
      paralelo ao eixo (perto dele) cruza o eixo no mesmo ponto, a uma distância
      $R/2$ do centro $C$ e, como $V = (R, 0)$, também a $R/2$ do vértice.
  \end{itemize}
\end{comentario}

A partir disso, como sabemos que o feixe que chega paralelo sempre vai cruzar o
eixo em $x_0 = R/2$, damos um nome para esse ponto, por ele ser tão importante:
chamamos de \textbf{foco} ($F$), e a distância $f = R/2$ do foco ao vértice, de
distância focal (Figura~\ref{fig:foco}).

\begin{figure}[htbp]
  \centering
  % Figura 7: feixes paralelos e o foco. C na origem, R = 6, F = (R/2, 0).
% Reflexão exata: o feixe na altura h cruza o eixo em x = R/(2 cos theta),
% sin(theta) = h/R; perto do eixo isso tende a R/2.
\begin{tikzpicture}[font=\small]
  \def\R{6}
  \coordinate (C) at (0,0);
  \coordinate (F) at (\R/2,0);
  \coordinate (V) at (\R,0);
  \draw[eixo] (-0.4,0) -- (6.6,0);
  \draw[espelho] (-26:\R) arc[start angle=-26,end angle=26,radius=\R];
  \foreach \a in {-25.5,-23,...,24.5}
    \draw[hachura] (\a:\R) -- ++(\a-40:0.18);
  \foreach \h/\c in {1.8/azul,1.2/laranja,0.6/verde,-0.6/rosa,-1.2/vermelhao,-1.8/vermelho} {
    \pgfmathsetmacro{\th}{asin(\h/\R)}
    \pgfmathsetmacro{\xM}{sqrt(\R*\R-\h*\h)}
    \pgfmathsetmacro{\xX}{\R/(2*cos(\th))}
    \coordinate (M) at (\xM,\h);
    \coordinate (X) at (\xX,0);
    \draw[raio,\c] (0.4,\h) -- (M);
    \draw[raio,\c] (M) -- ($(M)!1.45!(X)$);
  }
  \fill (C) circle (0.053);
  \fill (F) circle (0.07);
  \node[fill=white,inner sep=1pt] at (3.0,0.282) {$F$};
  \node[inner sep=0pt,anchor=base west] at (-0.411,-0.396) {$C$};
  \node[inner sep=0pt,anchor=base west] at (6.156,-0.396) {$V$};
  \foreach \x in {0,3,6} \draw[guia] (\x,0) -- (\x,-3.1);
  \draw[cota] (0,-2.95) -- (3,-2.95) node[midway,fill=white,inner sep=1.5pt] {$R/2$};
  \draw[cota] (3,-2.95) -- (6,-2.95) node[midway,fill=white,inner sep=1.5pt] {$f = R/2$};
\end{tikzpicture}

  \caption{Feixes que chegam paralelos ao eixo, em alturas diferentes, são
    refletidos pelo espelho e cruzam o eixo no mesmo ponto, a $R/2$ do centro e
    do vértice: esse ponto é o foco $F$. (Os feixes mais afastados do eixo cruzam
    um pouco antes, porque para eles $h \ll R$ vale menos; por isso os espelhos
    reais têm abertura pequena.)}
  \label{fig:foco}
\end{figure}

Note que, como o procedimento tomado foi rigoroso analiticamente, foi possível
mensurar de forma precisa o tamanho do erro das aproximações, assim garantindo a
validade da expressão. Denota-se tal procedimento por \textbf{aproximação
paraxial}. Utilizando, portanto, tal resultado, equacionemos os feixes refletidos
(Figura~\ref{fig:gauss}).

\begin{figure}[htbp]
  \centering
  % Figura 8: construção para a equação de Gauss (dimensões exageradas).
% C na origem; R = 2f = 7; O = (R - p, h) com p = 5,6 e h = 1,3.
% delta: chega paralelo, reflete em P e passa por F; gamma: reflete em V com o
% mesmo ângulo alpha (passa por (R - p, -h)). I = delta \cap gamma.
\begin{tikzpicture}[font=\small,scale=0.9]
  \def\R{7}
  \def\h{1.3}
  \def\p{5.6}
  \pgfmathsetmacro{\xP}{sqrt(\R*\R-\h*\h)}
  \pgfmathsetmacro{\al}{atan(\h/\p)}
  \coordinate (O) at ({\R-\p},\h);
  \coordinate (XO) at ({\R-\p},0);
  \coordinate (Ol) at ({\R-\p},-\h);
  \coordinate (P) at (\xP,\h);
  \coordinate (F) at ({\R/2},0);
  \coordinate (V) at (\R,0);
  \coordinate (I) at (intersection of P--F and V--Ol);
  % eixos
  \draw[line width=0.5pt,-{Stealth[length=4pt,width=3pt]}] (-2.6,0) -- (8.4,0);
  \draw[line width=0.5pt,-{Stealth[length=4pt,width=3pt]}] (0,-2.75) -- (0,2.45);
  \node[inner sep=0pt,anchor=base west] at (8.1,-0.29) {$x$};
  \node[inner sep=0pt,anchor=base west] at (-0.32,2.2) {$y$};
  % espelho
  \draw[espelho] (-16:\R) arc[start angle=-16,end angle=16,radius=\R];
  \foreach \a in {-15,-13,...,15}
    \draw[hachura] (\a:\R) -- ++(\a-40:0.2);
  % O e sua projeção X_O; cota p'
  \draw[line width=0.5pt,dashed] (O) -- (XO);
  \draw[guia] (I) -- ++(0,-0.62);
  \draw[guia] (V) -- ++(0,-2.67);
  \draw[cota] ($(I)+(0,-0.5)$) -- ($(V|-I)+(0,-0.5)$)
    node[midway,fill=white,inner sep=1.5pt] {$p'$};
  % feixe delta: chega paralelo, reflete em P e passa pelo foco
  \draw[raio,vermelho] (O) -- (P);
  \draw[raio,vermelho] (P) -- (F);
  \draw[raio,vermelho] (F) -- (I);
  \node[inner sep=0pt,anchor=base west,text=vermelho] at (4.0,1.43) {$\delta$};
  % feixe O -> V e seu reflexo gamma
  \draw[raio,azul] (O) -- (V);
  \draw[raio,azul] (V) -- (I);
  \node[inner sep=0pt,anchor=base west,text=azul] at (2.9,0.58) {$\gamma$};
  % ângulos
  \draw[line width=0.4pt] ($(F)+(0.9,0)$) arc[start angle=0,end angle={atan(\h/(\xP-\R/2))},radius=0.9];
  \node[inner sep=0pt,anchor=base west] at ($(F)+(1.02,0.07)$) {$2\theta$};
  \draw[line width=0.4pt,azul] ($(V)+(180-\al:1.1)$) arc[start angle=180-\al,end angle=180,radius=1.1];
  \node[inner sep=0pt,anchor=base west,text=azul] at ($(V)+(-1.48,0.07)$) {$\alpha$};
  % pontos e rótulos
  \fill (O) circle (0.07) (XO) circle (0.055) (P) circle (0.07) (F) circle (0.055)
    (V) circle (0.07) (I) circle (0.07);
  \node[inner sep=0pt,anchor=base west] at (0.5,1.52) {$O\,(R-p,\ h)$};
  \node[inner sep=0pt,anchor=base west] at (7.12,1.4) {$P$};
  \node[anchor=north,inner sep=2pt] at (XO) {$X_O$};
  \node[anchor=north west,inner sep=0pt] at ($(F)+(0.04,-0.1)$) {$F$};
  \node[inner sep=0pt,anchor=base west] at (6.68,-0.44) {$V$};
  \node[anchor=east,inner sep=3pt] at (I) {$I\,(2f-p',\ y_{p'})$};
\end{tikzpicture}

  \caption{Os feixes refletidos $\delta$ e $\gamma$, que saem do ponto $O$ do
    objeto, se cruzam no ponto $I$ da imagem (dimensões exageradas).}
  \label{fig:gauss}
\end{figure}

\noindent
\begin{minipage}[t]{0.49\textwidth}
  Para $\tg 2\theta$, considere $\triangle PVF$, com $\overline{FV} \approx f$:
  \[
    % Triângulo PVF (para tg 2theta), esquemático.
\begin{tikzpicture}[font=\small,baseline=(b)]
  \coordinate (b) at (0,0.4);
  \draw[line width=0.6pt] (0,0) -- (2.0,0) -- (2.0,1.0) -- cycle;
  \draw[line width=0.4pt] (2.0,0) ++(-0.16,0) -- ++(0,0.16) -- ++(0.16,0);
  \draw[line width=0.4pt] (0.6,0) arc[start angle=0,end angle={atan(1.0/2.0)},radius=0.6];
  \node[inner sep=0pt,anchor=base west] at (0.72,0.07) {$2\theta$};
  \node[anchor=north,inner sep=2pt] at (1.0,0) {$f$};
  \node[anchor=west,inner sep=2pt] at (2.0,0.5) {$h$};
\end{tikzpicture}
\qquad
    \tg 2\theta = m_\delta \approx \frac{h}{f}
  \]
\end{minipage}\hfill
\begin{minipage}[t]{0.47\textwidth}
  Para $\tg\alpha$, $\triangle OVX_O$:
  \[
    % Triângulo O V X_O (para tg alpha), esquemático.
\begin{tikzpicture}[font=\small,baseline=(b)]
  \coordinate (b) at (0,0.35);
  \draw[line width=0.6pt] (0,0) -- (2.4,0) -- (0,0.85) -- cycle;
  \draw[line width=0.4pt] (0.16,0) -- ++(0,0.16) -- ++(-0.16,0);
  \draw[line width=0.4pt,azul] (1.7,0) arc[start angle=180,end angle={180-atan(0.85/2.4)},radius=0.7];
  \node[inner sep=0pt,anchor=base east,text=azul] at (1.6,0.06) {$\alpha$};
  \node[anchor=north,inner sep=2pt] at (1.2,0) {$p$};
\end{tikzpicture}
\qquad
    \tg\alpha = m_\gamma = \frac{h}{p}
  \]
\end{minipage}

\medskip
Logo:
\[
  \begin{dcases}
    \text{1º)}\ \ \delta\colon\ y - \frac{h}{f}\,x - b_\delta = 0; &
      \text{1.1)}\ P = (x_P,\ h) \in \delta,\ x_P \approx 2f\\[2pt]
    \text{2º)}\ \ \gamma\colon\ y - \frac{h}{p}\,x - b_\gamma = 0; &
      \text{2.1)}\ V = (2f,\ 0) \in \gamma.
  \end{dcases}
\]
\begin{align*}
  \text{1.1)} &\Longrightarrow\ b_\delta = h - \frac{h}{f}\cdot 2f = -h\\
  \text{2.1)} &\Longrightarrow\ b_\gamma = 0 - \frac{h}{p}\cdot 2f = -\frac{2hf}{p}
\end{align*}
Assim:
\[
  \begin{dcases}
    \delta\colon\ y - \frac{h}{f}\,x + h = 0\\
    \gamma\colon\ y - \frac{h}{p}\,x + \frac{2hf}{p} = 0
  \end{dcases}
\]
Como $\delta \cap \gamma = \{I\}$,
\begin{align}
  &\Longrightarrow\ \cancel{y_{p'}} - \frac{\cancel{h}}{f}(2f - p') + \cancel{h}
    = \cancel{y_{p'}} - \frac{\cancel{h}}{p}(2f - p') + \frac{2\cancel{h}f}{p}
    \notag\\
  &\Longrightarrow\ -\frac{2\cancel{f}}{\cancel{f}} + \frac{p'}{f} + 1
    = \cancel{-\frac{2f}{p}} + \frac{p'}{p} + \cancel{\frac{2f}{p}} \notag\\
  &\Longrightarrow\ -1 + \frac{p'}{f} = \frac{p'}{p} \notag\\
  &\Longrightarrow\ \left(\frac{1}{p'}\right)\left[-1 + \frac{p'}{f}\right]
    = \left(\frac{1}{p'}\right)\frac{p'}{p} \notag\\
  &\Longrightarrow\ \frac{1}{f} - \frac{1}{p'} = \frac{1}{p}
    \quad\Longrightarrow\quad
    \underset{\text{equação de Gauss}}{\boxed{\frac{1}{f} = \frac{1}{p} + \frac{1}{p'}}}
    \qquad \text{C.Q.D.} \label{eq:gauss}
\end{align}

\bigskip
\begin{center}
  \small\itshape (continua)
\end{center}

\end{document}
